% LaTeX code: Introduction through Background Knowledge only
% Same format as full thesis (dalthesis, 12pt, same packages)

\documentclass[12pt]{dalthesis}

\usepackage{geometry}
\geometry{a4paper, margin=1in}

\usepackage{longtable}
\usepackage{array}
\usepackage{nomencl}
\usepackage{booktabs}
\usepackage{multirow}
\usepackage{tabularx}
\usepackage{adjustbox}
\usepackage{graphicx}
\usepackage{caption}
\usepackage{float}
\usepackage{pdflscape}
\usepackage{amsmath}
\usepackage{amssymb}
\usepackage{algorithm}
\usepackage{algorithmic}
\usepackage{hyperref}
\usepackage{multicol}
\usepackage{placeins}
\usepackage{listings}
\usepackage{xcolor}
\usepackage{colortbl}
\usepackage{pifont}
\makenomenclature

\definecolor{codegreen}{rgb}{0,0.6,0}
\definecolor{codegray}{rgb}{0.5,0.5,0.5}
\lstset{
    basicstyle=\ttfamily\small,
    breaklines=true,
    frame=single,
    numbers=left,
    numberstyle=\tiny,
    keywordstyle=\color{blue},
    commentstyle=\color{codegreen},
}
\newcommand{\xmark}{\ding{55}}

\begin{document}

\title{Explainable AI for Diabetes Risk Prediction Using Grad-CAM and IGTD}
\author{Name of the student}

\macs  
\degreeinitial{M.A.C.Sc.}
\faculty{Computer Science}
\dept{Faculty of Computer Science}

\defencemonth{Dec}
\defenceyear{2024}

\makenomenclature
\renewcommand{\nomname}{List of Abbreviations and Symbols Used}

\nomenclature[A]{\textbf{CDC}}{Centers for Disease Control and Prevention}
\nomenclature[A]{\textbf{CNN}}{Convolutional Neural Network}
\nomenclature[A]{\textbf{ENN}}{Edited Nearest Neighbors}
\nomenclature[A]{\textbf{Grad-CAM}}{Gradient-weighted Class Activation Mapping}
\nomenclature[A]{\textbf{IGTD}}{Image Generator for Tabular Data}
\nomenclature[A]{\textbf{XAI}}{Explainable Artificial Intelligence}

\frontmatter

\begin{abstract}
This report presents an Explainable AI (XAI) approach for diabetes risk prediction using a Convolutional Neural Network (CNN) combined with the Image Generator for Tabular Data (IGTD) algorithm. We transform tabular health data from the CDC diabetes dataset into 2D grayscale images, train a CNN based on Shenghao Wang et al.'s cnn\_igtd\_f15\_enn model, and apply Grad-CAM to visualize which features the model focuses on when making predictions. The pipeline achieves 84\% accuracy on the test set and produces interpretable heatmaps with feature names and color scales, supporting trust and transparency in the model's decisions. The implementation is available on GitHub: \url{https://github.com/mzaman2016/VitaScreen}.
\end{abstract}

\addcontentsline{toc}{chapter}{List of Abbreviations and Symbols Used}
\printnomenclature

\begin{acknowledgements}
% Add your acknowledgements here.
\end{acknowledgements}

\mainmatter

\chapter{Introduction}
\section{Overview}

Diabetes is a major global health concern, and early risk assessment can support preventive care. Machine learning models can predict diabetes risk from health indicators, but black-box models lack interpretability. This work applies Explainable AI (XAI) using Grad-CAM to a CNN that predicts diabetes from tabular data transformed into images via IGTD.

\section{Problem Statement}

Deep learning models for healthcare often lack transparency. Clinicians and patients need to understand \textit{why} a model predicts a certain outcome. We address this by integrating Grad-CAM into a diabetes prediction pipeline to highlight which input features drive each prediction.

\section{Motivation}

Trust in AI systems requires explainability. Grad-CAM provides visual explanations by showing which regions of the input (IGTD images) the CNN focuses on. This helps validate model behavior and identify potential biases.

\section{Major Contributions}

\begin{itemize}
\item Implemented a full pipeline: CDC dataset $\rightarrow$ IGTD $\rightarrow$ CNN $\rightarrow$ Grad-CAM.
\item Added feature names and color scales to Grad-CAM visualizations for interpretability.
\item Achieved 84\% accuracy on the test set with the cnn\_igtd\_f15\_enn architecture.
\item Integrated with the VitaScreen GitHub repository.
\end{itemize}

\section{Organization of the report}

Chapter~\ref{ch:background} covers background on IGTD, CNNs, and Grad-CAM. Chapter~\ref{ch:lit} reviews related work. Chapter~\ref{ch:method} describes the methodology. Chapter~\ref{ch:impl} presents implementation details. Chapter~\ref{ch:results} discusses results. Chapter~\ref{ch:conclusion} concludes with limitations and future work.

\chapter{Background Knowledge}\label{ch:background}

\section{IGTD (Image Generator for Tabular Data)}

IGTD converts tabular data into 2D images by optimizing the arrangement of features in a grid so that similar features are placed close together. We use a 5$\times$3 grid for 15 features.

\section{CNN Architecture}

We use Shenghao Wang et al.'s cnn\_igtd\_f15\_enn: 4 convolutional blocks (64 filters, 3$\times$3), BatchNorm, ReLU, Global Average Pooling, Dropout (50\%), and a fully connected output layer.

\section{Grad-CAM}

Grad-CAM computes a heatmap by combining the gradients of the target class with respect to the last convolutional layer's activations. Red indicates high importance; blue indicates low importance.

\section{ENN (Edited Nearest Neighbors)}

ENN removes noisy or borderline samples from the training set using k-Nearest Neighbors (k=3). We apply ENN only on the training set before IGTD.

\section{CDC Diabetes Dataset}

The CDC diabetes dataset contains health indicators: HighBP, HighChol, CholCheck, BMI, Smoker, Stroke, HeartDiseaseorAttack, PhysActivity, Fruits, Veggies, HvyAlcoholConsump, AnyHealthcare, NoDocbcCost, GenHlth, MentHlth. We use 15 features and a binary label: 0 = Non-Diabetic, 1 = Diabetic (or prediabetic).

\chapter{Literature Review}\label{ch:lit}

\textbf{What I will do:} I will read and summarize papers about IGTD (Zhu et al.), Grad-CAM (Selvaraju et al.), and diabetes prediction with CNNs (Wang et al., VitaScreen). I will compare our work with these papers and explain how our approach is similar or different.

% TODO: Add your literature review content here.

\chapter{Proposed Methodology}\label{ch:method}

\textbf{What I will do:} I will describe the full pipeline step by step: (1) Load CDC data, (2) Select 15 features, (3) Split 80:20 train/test, (4) Apply ENN on the training set, (5) Run IGTD to create images, (6) Train the CNN, (7) Apply Grad-CAM and save heatmaps with feature names and scales. I will also explain how we add feature names and color scales to the Grad-CAM images.

% TODO: Add your methodology content here.

\chapter{Implementation}\label{ch:impl}

\textbf{What I will do:} I will explain the code. I will describe the base pipeline (diabetes\_prediction\_pipeline.py) and the XAI pipeline (gradcam\_explainability.py). I will show the main functions and how they connect. I will also give the GitHub link and how to run the code.

% TODO: Add your implementation content here.

\chapter{Result Discussion}\label{ch:results}

\textbf{What I will do:} I will show the results table (Accuracy 0.84, Precision 0.38, Recall 0.50, F1-Score 0.43). I will add the 8 Grad-CAM sample figures. I will discuss what the heatmaps show—for example, that Veggies, AnyHealthcare, and CholCheck are often important. I will also explain why the model does better on non-diabetic cases (because there are more of them in the data).

% TODO: Add your results, tables, and figures here.

\chapter{Conclusion}\label{ch:conclusion}

\textbf{What I will do:} I will summarize the whole project. I will say that we built an explainable diabetes prediction pipeline using IGTD and Grad-CAM, achieved 84\% accuracy, and added feature names and scales so people can understand the model's decisions. I will mention the GitHub link again.

% TODO: Add your conclusion content here.

\chapter{Limitations and Future Work}

\section{Limitations}

\textbf{What I will do:} I will say that (1) there are more non-diabetic than diabetic samples, so the model is better at non-diabetic cases; (2) Grad-CAM shows image regions, not direct feature importance; (3) we used 500 samples, so results may change with more data.

% TODO: Add your limitations content here.

\section{Future Work}

\textbf{What I will do:} I will say we could (1) use SHAP or LIME for feature-level explanations; (2) use the full CDC dataset or the cdcNormalDiabeticFE1\_20RFFSQ.csv file; (3) put Grad-CAM into a web app; (4) try NCTD instead of IGTD.

% TODO: Add your future work content here.

\end{document}
